\documentclass{abntex2}
\usepackage[utf8]{inputenc}
\usepackage[T1]{fontenc}
\usepackage{geometry}
\usepackage{graphicx}
\usepackage{float}
\usepackage{hyperref}
\usepackage[english,indonesian]{babel}

\title{Optimasi Pemrosesan Template LaTeX Menggunakan Pendekatan Berbasis AI}
\author{Academic Support Agent}
\date{\today}

\begin{document}
\maketitle

\begin{abstract}

\end{abstract}

\chapter{PENDAHULUAN}


\chapter{TINJAUAN PUSTAKA}


% Original content (commented out):
%     %% abtex2-modelo-trabalho-academico.tex, v-1.7.1 laurocesar
    %% Copyright 2012-2013 by abnTeX2 group at http://abntex2.googlecode.com/ 
    %%
    %% This work may be distributed and/or modified under the
    %% conditions of the LaTeX Project Public License, either version 1.3
    %% of this license or (at your option) any later version.
    %% The latest version of this license is in
    %%   http://www.latex-project.org/lppl.txt
    %% and version 1.3 or later is part of all distributions of LaTeX
    %% version 2005/12/01 or later.
    %%
    %% This work has the LPPL maintenance status `maintained'.
    %% 
    %% The Current Maintainer of this work is the abnTeX2 team, led
    %% by Lauro César Araujo. Further information are available on 
    %% http://abntex2.googlecode.com/
    %%
    %% This work consists of the files abntex2-modelo-trabalho-academico.tex,
    %% abntex2-modelo-include-comandos and abntex2-modelo-references.bib
    %%

    % ------------------------------------------------------------------------
    % ------------------------------------------------------------------------
    % abnTeX2: Modelo de Trabalho Academico (tese de doutorado, dissertacao de
    % mestrado e trabalhos monograficos em geral) em conformidade com 
    % ABNT NBR 14724:2011: Informacao e documentacao - Trabalhos academicos -
    % Apresentacao
    % ------------------------------------------------------------------------
    % ------------------------------------------------------------------------

    \documentclass[
	    % -- opções da classe memoir --
	    12pt,				% tamanho da fonte
	    openright,			% capítulos começam em pág ímpar (insere página vazia caso preciso)
	    oneside,			% para impressão em verso e anverso. Oposto a oneside 'twoside'
	    a4paper,			% tamanho do papel. 
	    % -- opções da classe abntex2 --
        %chapter=TITLE,		% títulos de capítulos convertidos em letras maiúsculas
	    %section=TITLE,		% títulos de seções convertidos em letras maiúsculas
	    %subsection=TITLE,	% títulos de subseções convertidos em letras maiúsculas
	    %subsubsection=TITLE,% títulos de subsubseções convertidos em letras maiúsculas
	    % -- opções do pacote babel --
	english,			% idioma adicional para hifenização
	french,				% idioma adicional para hifenização
	spanish,			% idioma adicional para hifenização
	brazil,				% o último idioma é o principal do documento
	]{abntex2}


    % ---
    % PACOTES
    % ---
    \usepackage[brazilian,hyperpageref]{backref}	 % Paginas com as citações na bibl
    \usepackage[alf]{abntex2cite}	% Citações padrão ABNT

    \usepackage{float} % Melhorar as distâncias das imagens no cabeçalho

    % ---
    % Pacotes fundamentais 
    % ---
    \usepackage{cmap}				% Mapear caracteres especiais no PDF
    \usepackage{lmodern}			% Usa a fonte Latin Modern			
    \usepackage[T1]{fontenc}		% Selecao de codigos de fonte.
    \usepackage[utf8]{inputenc}		% Codificacao do documento (conversão automática dos acentos)
    \usepackage{lastpage}			% Usado pela Ficha catalográfica 
    \usepackage{indentfirst}		% Indenta o primeiro parágrafo de cada seção.
    \usepackage{color}				% Controle das cores
    \usepackage{graphicx}			% Inclusão de gráficos
    %\usepackage{verbatim}           % Comentário de multiplas linhas
    \usepackage{enumerate}          % Enumerar em difernetes estilos A, a, I, i ou 1
    \usepackage[table,xcdraw]{xcolor}
    \usepackage{amssymb}            % Símbolo matemático
    %\usepackage{pdflscape}
    \usepackage{amsmath}
    \usepackage{subfig}
    \usepackage{multicol} 
    %\usepackage{tablefootnote}
    %\usepackage{ulem} % sublinhar, estava aparecendo nas referências






    %\usepackage{floatrow}
    % Table float box with bottom caption, box width adjusted to content
    %\newfloatcommand{capbtabbox}{table}[][\FBwidth]

    \usepackage[tableposition=top]{caption}  % Eu inseri caption de tabela no topo
    %\usepackage[figureposition=top]{caption}

    \usepackage{blindtext}

   % \usepackage[utf8]{inputenc}
    \usepackage{geometry}
    \usepackage{algorithm}
    \usepackage{algpseudocode}

    % Pseudo Código
 %   \usepackage{algpseudocode,algorithm} 
    % Declaracoes em Português
    \algrenewcommand\algorithmicend{\textbf{end}}
    \algrenewcommand\algorithmicdo{\textbf{faça}}
    \algrenewcommand\algorithmicwhile{\textbf{ENQUANTO}}
    \algrenewcommand\algorithmicfor{\textbf{for}}
    \algrenewcommand\algorithmicif{\textbf{SE}}
    \algrenewcommand\algorithmicthen{\textbf{ENTÃO}}
    \algrenewcommand\algorithmicelse{\textbf{SENÃO}}
    \algrenewcommand\algorithmicreturn{\textbf{SAÍDA}}
    \algrenewcommand\algorithmicfunction{\textbf{FUNÇÃO}}

    % Rearranja os finais de cada estrutura
    \algrenewtext{EndWhile}{\algorithmicend\ \algorithmicwhile}
    \algrenewtext{EndFor}{\algorithmicend\ \algorithmicfor}
    \algrenewtext{EndIf}{\algorithmicend\ \algorithmicif}
    \algrenewtext{EndFunction}{\algorithmicend\ \algorithmicfunction}

    % O comando For, a seguir, retorna 'para #1 -- #2 até #3 faça'
    \algnewcommand\algorithmicto{\textbf{até}}
    \algrenewtext{For}[3]%
    {\algorithmicfor\ #1 $\gets$ #2 \algorithmicto\ #3 \algorithmicdo}
    % ---
    %%
    \usepackage{cancel} % Insere zero na expressão.
    % ---
    % Pacotes adicionais, usados apenas no âmbito do Modelo Canônico do abnteX2
    % ---
    \usepackage{lipsum}				% para geração de dummy text
    % ---
    % Define o caminho das figuras
    \graphicspath{{figuras/}}
    % ---
    % Pacotes de citações
    % ---
    \usepackage[brazilian,hyperpageref]{backref}	 % Paginas com as citações na bibl
    \usepackage[alf]{abntex2cite}	% Citações padrão ABNT


    % --- 
    % CONFIGURAÇÕES DE PACOTES
    % --- 

    % ---
    %\begin{comment}
    % Configurações do pacote backref
    % Usado sem a opção hyperpageref de backref
    \renewcommand{\backrefpagesname}{Citado na(s) página(s):~}
    % Texto padrão antes do número das páginas
    \renewcommand{\backref}{}
    % Define os textos da citação
    \renewcommand*{\backrefalt}[4]{
	\ifcase #1 %
		Nenhuma citação no texto.%
	\or
		Citado na página #2.%
	\else
		Citado #1 vezes nas páginas #2.%
	\fi}%




    % ---
    %end{comment}

    % ---
    % Informações de dados para CAPA e FOLHA DE ROSTO
    % ---
 
   
    \titulo{TÍTULO} 

    \autor{\large Nome do Autor}
    \local{São Cristóvão -- SE, Brasil}
    \data{Janeiro de 2019}
    \orientador{Prof.Dr. Nome do Orientador}
    %\coorientador{xxxx}
    %\instituicao{%
    %  UNIVERSIDADE FEDERAL DE SERGIPE
    %  \par
    %  CENTRO DE CIÊNCIAS EXATAS E TECNOLOGIA
    %  \par
    %  PROGRAMA DE PÓS-GRADUAÇÃO EM ENGENHARIA ELÉTRICA}
    \tipotrabalho{Dissertação (Mestrado)}
    % O preambulo deve conter o tipo do trabalho, o objetivo, 
    % o nome da instituição e a área de concentração 
    \preambulo{Dissertação de Mestrado apresentada ao Programa de Pós-graduação em Engenharia Elétrica – PROEE, da Universidade Federal de Sergipe, como parte dos
    requisitos necessários à obtenção do título de Mestre em
    Engenharia Elétrica}
    % ---


    % ---
    % Configurações de aparência do PDF final

    % alterando o aspecto da cor azul
    \definecolor{blue}{RGB}{45,5,145} % anterior {41,5,145}
    \definecolor{red}{RGB}{255,0,0} % anterior {41,5,145}

    % informações do PDF
    \makeatletter
    \hypersetup{
     	%pagebackref=true,
		pdftitle={\@title}, 
		pdfauthor={\@author},
    	pdfsubject={\imprimirpreambulo},
	    pdfcreator={LaTeX with abnTeX2},
		pdfkeywords={abnt}{latex}{abntex}{abntex2}{trabalho acadêmico}, 
		colorlinks=true,       		% false: boxed links; true: colored links
    	linkcolor=blue,          	% color of internal links
    	citecolor=blue,        		% color of links to bibliography
    	filecolor=magenta,      		% color of file links
		urlcolor=blue,
		bookmarksdepth=4
    }
    \makeatother
    % --- 

    % --- 
    % Espaçamentos entre linhas e parágrafos 
    % --- 

    % O tamanho do parágrafo é dado por:
    \setlength{\parindent}{1.3cm}

    % Controle do espaçamento entre um parágrafo e outro:
    \setlength{\parskip}{0.2cm}  % tente também \onelineskip

    % ---
    % compila o indice
    % ---
    \makeindex
    % ---

    % ----
    % Início do documento
    % ----
    \begin{document}

    % Retira espaço extra obsoleto entre as frases.
    \frenchspacing 

    % ----------------------------------------------------------
    % ELEMENTOS PRÉ-TEXTUAIS
    % ----------------------------------------------------------
    % \pretextual





    %\imprimircapa
    % ---
    % ---
    % Capa
    % ---
    \begin{figure}[!htb]
        \centering
        \includegraphics[scale=0.07]{figuras/LogoNovaUFS.png}
        %\includegraphics[scale=0.7]{figuras/ufs}
        % \includegraphics[width=3.6cm,height=3.6cm]{figuras/ufs}
    \end{figure}

    \begin{center}
        %\vspace*{0.0cm}
        \par
        \textbf{UNIVERSIDADE FEDERAL DE SERGIPE}
        \par
        \textbf{PRÓ-REITORIA DE PÓS-GRADUAÇÃO E PESQUISA}
        \par
        \textbf{PROGRAMA DE PÓS-GRADUAÇÃO EM ENGENHARIA ELÉTRICA}
        \vspace*{3cm}
        \par
        \textbf{\imprimirtitulo}
        %\textbf{DETECÇÃO DE FALHAS EM MATERIAIS COM TERMOGRAFIA ATIVA}
        %\vspace*{0.2cm}
        %   \par
        %   \textbf{\Large{ao \textit{Reality Gap}}}
        %\par
        %\textbf{Proposta de Dissertação}
        \vspace{2.5cm}
        \par
        \textbf{Nome do Autor}
        \vspace{2cm}
        \par
    %    Orientador: Prof.Dr. Nome do Orientador
        \par
    %    Co-orientador: xxx
        \vspace{1cm}
        {\imprimirlocal}
        \par
        {\imprimirdata}
        %\vspace{1cm}
    
  \end{center}
    \newpage


    % ---
    % Folha de rosto
    % (o * indica que haverá a ficha bibliográfica)

    % Capa
    % ---
    \begin{figure}[!htb]
    \centering
    \includegraphics[scale=0.07]{figuras/LogoNovaUFS.png}
    %\includegraphics[scale=0.7]{figuras/ufs}
   % \includegraphics[width=3.6cm,height=3.6cm]{figuras/ufs}
    \end{figure}

    \begin{center}
        %\vspace*{0.0cm}
        \par
        \vspace*{1cm}
        \par
        \textbf{\imprimirtitulo}
        %\textbf{DETECÇÃO DE FALHAS EM MATERIAIS COM TERMOGRAFIA ATIVA}
        \par
        \vspace*{1cm}
        \par
        \textbf{Nome do Autor}
        \vspace*{1cm}
    \end{center}


    \hspace{.3\textwidth}
    \begin{minipage}{.6\textwidth}
    \SingleSpacing
    \imprimirpreambulo
    \end{minipage}%
    %\vspace*{\fill}

    \hspace{.3\textwidth}
     \begin{minipage}{.6\textwidth}
    \SingleSpacing
     Orientador: Prof.Dr. Nome do Orientador
    \par
    %Co-orientador: xxxxxx
    \end{minipage}%
    \vspace*{\fill}


    \begin{center}
        %\vspace*{4cm}
        %\par
        %{São Cristóvão-SE, Brasil}
        %\par
        %{Fevereiro de 2018}
        {\imprimirlocal}
        \par
        {\imprimirdata}
        \vspace*{1cm}
    \end{center} 
    \newpage

% ---
    %\imprimirfolhaderosto* %(COMENTEI ESTA LINHA)
% ---

% ---
% Inserir a ficha bibliografica
% ---

% Isto é um exemplo de Ficha Catalográfica, ou ``Dados internacionais de
% catalogação-na-publicação''. Você pode utilizar este modelo como referência. 
% Porém, provavelmente a biblioteca da sua universidade lhe fornecerá um PDF
% com a ficha catalográfica definitiva após a defesa do trabalho. Quando estiver
% com o documento, salve-o como PDF no diretório do seu projeto e substitua todo
% o conteúdo de implementação deste arquivo pelo comando abaixo:
%
% \begin{fichacatalografica}
%     \includepdf{fig_ficha_catalografica.pdf}
% \end{fichacatalografica}

% Antigo formato da folha de aprovação

% ---
% Inserir folha de aprovação
% ---

% Isto é um exemplo de Folha de aprovação, elemento obrigatório da NBR
% 14724/2011 (seção 4.2.1.3). Você pode utilizar este modelo até a aprovação
% do trabalho. Após isso, substitua todo o conteúdo deste arquivo por uma
% imagem da página assinada pela banca com o comando abaixo:
%
% \includepdf{folhadeaprovacao_final.pdf}
%
\begin{folhadeaprovacao}

  \begin{center}
  
  
    %{\ABNTEXchapterfont\large\imprimirautor}
    
    %\vspace*{\fill}\vspace*{\fill}
    {\ABNTEXchapterfont\bfseries\imprimirtitulo}
    \vspace*{1.5cm}
    
    {\ABNTEXchapterfont\large{\textbf{Nome do Autor}}}
    \vspace*{1cm}

  %  \hspace{.45\textwidth}
  %  \begin{minipage}{.5\textwidth}
  %     \imprimirpreambulo
  %  \end{minipage}%
  %  \vspace*{\fill}
  \end{center}
  
\noindent \textbf{Dissertação de Mestrado apresentada ao Programa de Pós-graduação em Engenharia Elétrica – PROEE, da
Universidade Federal de Sergipe, como parte dos requisitos necessários à obtenção do título de Mestre em Engenharia Elétrica} \\


\noindent Examinada por:

   %\assinatura{\textbf{\imprimirorientador} \\ Orientador} 
   \assinatura{\textbf{Prof. Dr. A... \\ DEMEC/UFS} \\ }
   %\vspace*{0.5cm}
   \assinatura{\textbf{Prof. Dr. C... \\P2CEM/ UFS} \\ }
   %\vspace*{0.5cm}
   \assinatura{\textbf{Prof. Dr. E... \\ PROEE/UFS} \\ }
   %\assinatura{\textbf{Professor} \\ Convidado 4}
    
   % \vspace*{\fill}  \vspace*{\fill}
   \begin{center}
    \vspace{1.0cm}
    {\large\imprimirlocal}
    \par
    {\large{30 de janeiro de 2019}}
    %\vspace*{1cm}
  \end{center}
  
\end{folhadeaprovacao}
% ---

\begin{comment}


% ---
% Dedicatória
% ---
\begin{dedicatoria}
   \vspace*{\fill}
   \centering
   \noindent
   \textit{ Dedicatória} \vspace*{\fill}
\end{dedicatoria}
% ---
\end{comment}



% ---
% Agradecimentos
% ---
    \begin{agradecimentos}

    Primeiramente agradeço ...


    \end{agradecimentos}
    % ---


   % \begin{comment}


    % ---
 %   Epígrafe
    % ---
    \begin{epigrafe}
        \vspace*{\fill}
	    \begin{flushright}
	    "O estudo, a busca da verdade e da beleza são domínios em que nos é consentido sermos crianças por toda a vida."\\
	    \vspace*{0.5cm}
                                                Albert Einstein
	    \end{flushright}
    \end{epigrafe}
    % ---

   % \end{comment}

    % RESUMO WELLINGTON
    % ---
    \noindent Resumo da Dissertação apresentada ao PROEE/UFS como parte dos requisitos necessários para a obtenção do grau de Mestre (Me.)

    \begin{center}
        %\vspace*{0.0cm}
        \par
        \vspace*{0.15cm}
        \par
        \textbf{\imprimirtitulo}
        %\textbf{DETECÇÃO DE FALHAS EM MATERIAIS ISOLANTES COM TERMOGRAFIA ATIVA PULSADA}
        \par
        \vspace*{0.15cm}
        \par
        \textbf{Nome do Autor}
        \vspace*{0.15cm}
        \par
        \textbf{Janeiro/2019}
        \vspace*{0.15cm}
    \end{center}
        \par
    \noindent Orientador: Prof. Dr. Nome do Orientador
        \par
    %\noindent Co-orientador: xxxx
    \vspace*{0.15cm}
        \par

        Esta dissertação de mestrado apresenta o desenvolvimento de um algoritmo capaz de detectar regiões defeituosas em materiais utilizando como base a técnica de reconstrução de sinais termográficos e os princípios físicos de transferência de calor. Como base experimental, utilizou-se ensaios realizados e cedidos pelo Laboratório de Ensaio Não Destrutivo, Corrosão e Soldagem (LNDC) da UFRJ. Experimentalmente, utiliza-se um curto pulso térmico e monitora-se a temperatura da superfície do objeto ao longo do tempo através de uma câmera termográfica. Devido ao volume de dados, a primeira etapa do algoritmo é a compactação. Utilizou-se como base do fenômeno a lei de resfriamento de Newton e obteve-se uma taxa de compactação de 99\%.  As principais contribuições do algoritmo desenvolvido são: compactação e tratamento dos dados de aquecimento (que são negligenciados pelas outras técnicas); utilização do conceito de mudança no sentido do fluxo de calor para delimitar as bordas das regiões de interesse (onde provavelmente estão os defeitos); detecção das bordas/interface entre o objeto e o fundo da imagem; consolidação em uma única imagem, agregando os indicadores referentes ao conceito de constante de tempo de resfriamento/aquecimento, de máxima temperatura de cada \emph{pixel} e do contraste máximo. Finalmente, como resultado, apresenta-se uma única imagem com todas essas informações, diminuindo a subjetividade do operador e auxiliando na eficácia dessa técnica termográfica. 
    
    \par
    \noindent 
   \textbf{Palavras-chaves}: Termografia Ativa, contraste máximo, constante tempo, compactação, máximos temperatura.
    \newpage
    \newpage

    % Abstract Nome do Autor
    % ---
    \noindent Abstract of Dissertation presented to PROEE/UFS as a partial fulfillment of the requirements for the degree of Master

    \begin{center}
         %\vspace*{0.0cm}
        \par
        \vspace*{0.2cm}
        \par
        \textbf{DEVELOPMENT OF AN ALGORITHM FOR PROCESSING RADIOMETRIC DATA OBTAINED BY PULSED ACTIVE THERMOGRAPHY}
        \par
        \vspace*{0.2cm}
        \par
        \textbf{Nome do Autor}
        \vspace*{0.2cm}
        \par
        \textbf{Janeiro/2019}
        \vspace*{0.2cm}
    \end{center}
    \par
    \noindent Advisor: Prof. Dr. Nome do Orientador
    \par
    %\noindent Co-advisor: xxxxx
    \vspace*{0.2cm}
    \par
    This master thesis presents the development of an algorithm capable of detecting defects in materials, based on the technique of reconstruction of thermographic signals and the physical principles of heat transfer. As an experimental basis, tests performed and assigned by the Non-Destructive Testing, Corrosion and Welding Laboratory (LNDC) of UFRJ were used. Experimentally, a short thermal pulse is used and the surface temperature of the object is monitored over time through a thermographic camera. Due to the volume of data, the first step of the algorithm is data compression. Newton's law of cooling was used as the basis of the phenomenon and a compression ratio of 99 \% was obtained. The main contributions of the developed algorithm are: data compression and treatment of the heating data (which are neglected by the other techniques); use of the concept of change in the direction of the heat flow to delimit the edges of the regions of interest (Where the defects are probably located); edge detection/interface between the object and the image background; consolidation in a single image by aggregating the indicators referring to the concept of cooling/heating time constant, maximum temperature of each \emph{pixel} and maximum contrast. Finally, as a result, a single image with all this information is presented, reducing the subjectivity of the operator and improving the effectiveness of this thermographic technique.
    \par
    \noindent 
   \textbf{Keywords}: Active Thermography, maximum contrast, time constant, data compression, maximum temperature.
    \newpage








    \pagenumbering{Roman} % inicia numeração arábica

    % ------------------------------------------------------------------------
    % inserir lista de ilustrações
    % ------------------------------------------------------------------------
    \pdfbookmark[0]{\listfigurename}{lof}
    \listoffigures
    \cleardoublepage

    % ---
    % ------------------------------------------------------------------------
    % inserir lista de tabelas
    % ------------------------------------------------------------------------
    \pdfbookmark[0]{\listtablename}{lot}
    \listoftables
    \cleardoublepage
% ---
% ------------------------------------------------------------------------
% inserir lista de abreviaturas e siglas
% ------------------------------------------------------------------------
%\begin{siglas}
%  \item[Fig.] Area of the $i^{th}$ component
%  \item[456] Isto é um número
%  \item[123] Isto é outro número
%  \item[lauro cesar] este é o meu nome
%\end{siglas}
% ---
% ------------------------------------------------------------------------
% inserir lista de símbolos
% ------------------------------------------------------------------------
%\begin{simbolos}
%  \item[$ \Gamma $] Letra grega Gama
%  \item[$ \Lambda $] Lambda
%  \item[$ \zeta $] Letra grega minúscula zeta
%  \item[$ \in $] Pertence
%\end{simbolos}
% ---
% ------------------------------------------------------------------------
% inserir o sumario
% ------------------------------------------------------------------------

\pdfbookmark[0]{\contentsname}{toc}
\tableofcontents*
%\listoftables
\thispagestyle{empty} % não mostra numeração nesta página

%\cleardoublepage
% ---
% ------------------------------------------------------------------------
% ELEMENTOS TEXTUAIS
% ------------------------------------------------------------------------

\textual


% ========================================================================
% Introdução
% ========================================================================

\pagenumbering{arabic} % inicia numeração normal
%\include{introducao}

% ========================================================================
% Fundamentação Teórica
% ========================================================================

%\include{fundamentacao_teorica}

% ========================================================================
% Metodologia
% ========================================================================

%\include{desenvolvimento_algoritmo}

% ========================================================================
% Cronograma
% ========================================================================

%\include{cronograma}

% ========================================================================
% Resultados
% ========================================================================

%\include{resultados}


% ========================================================================
% Conclusão
% ========================================================================

%\include{conclusao}

% ----------------------------------------------------------
% Referências bibliográficas
% ----------------------------------------------------------
\bibliography{referencias}

%\include{apendice}


\end{document}


\end{document}